\SourcePage{330}{335}
\section{Electron propagator}

The concept of propagators (Green's functions) introduced in the preceding sections plays a fundamental role in the apparatus of quantum electrodynamics. The photon propagator $D_{\mu\nu}$, occupying in it a position as the basic quantity characterising the interaction of two electrons, is the same role that the electron propagator plays in the interaction of the electron and the photon.

Let us now turn to the actual computation of the propagators, starting with the electronic case. We act on the function
\begin{equation}
G_{ik}(x - x') = -i\langle 0|\mathrm{T}\psi_i(x)\bar{\psi}_k(x')|0\rangle
\tag{75.1}
\label{eq:75.1}
\end{equation}
($i$, $k$ are bispinor indices, where $\hat{p}_\mu = i\partial_\mu$). Since the operator $\hat{\psi}(x)$ satisfies the Dirac equation $(\gamma\hat{p} - m)\hat{\psi}(x) = 0$, we obtain zero at all points $x$, with the exception only of those for which $t = t'$. The point is that $G(x - x')$ tends to different limits when $t \to t' + 0$ and $t \to t' - 0$: according to definition~(74.8) these limits are respectively
\[
-i\langle 0|\psi_i(\mathbf{r},t)\bar{\psi}_k(\mathbf{r}',t)|0\rangle\quad\text{and}\quad +i\langle 0|\bar{\psi}_k(\mathbf{r}',t)\psi_i(\mathbf{r},t)|0\rangle
\]
and, as we shall see, on the light cone these do not coincide. This leads to the appearance in the derivative $\partial G/\partial t$ of an additional term with a $\delta$-function:
\begin{equation}
\frac{\partial G}{\partial t} = -i\langle 0|\mathrm{T}\frac{\partial\psi_i(x)}{\partial t}\bar{\psi}_k(x')|0\rangle + \delta(t - t')(G|_{t\to t'+0} - G|_{t\to t'-0}).
\tag{75.2}
\end{equation}

Noting that the operator $\gamma\hat{p} - m$ produces the derivative with respect to $t$ in the form $i\gamma^0\partial/\partial t$, we have therefore
\begin{equation}
(\gamma\hat{p} - m)_{ik}\,G_{kl}(x - x') = \delta(t - t')\gamma^0_{ik}\langle 0|\{\psi_k(\mathbf{r},t),\bar{\psi}_l(\mathbf{r}',t)\}_+|0\rangle.
\tag{75.3}
\end{equation}

We compute the anticommutator standing here. Multiplying the operators $\hat{\psi}(\mathbf{r},t)$ and $\hat{\bar{\psi}}(\mathbf{r}',t)$ (see~(73.6)) and using the commutation rules for fermionic operators $\hat{a}_\mathbf{p}$, $\hat{b}_\mathbf{p}$, we find
\begin{equation}
\{\hat{\psi}_i(\mathbf{r},t),\hat{\bar{\psi}}_k(\mathbf{r}',t)\}_+ = \sum_\mathbf{p}[\psi_{pi}(\mathbf{r})\bar{\psi}_{pk}(\mathbf{r}') + \psi_{-pi}(\mathbf{r})\bar{\psi}_{-pk}(\mathbf{r}')],
\tag{75.4}
\end{equation}

\SourcePage{331}{336}
where $\psi_{\pm\mathbf{p}}(\mathbf{r})$ are the wave functions without the time factor (as in~\S\S\,73, 74, for brevity we do not write out the polarisation indices). The totality of all functions $\psi_{\pm\mathbf{p}}(\mathbf{r})$---the eigenfunctions of the Hamiltonian of the electron---constitutes a complete system of normalised functions, and according to the general properties of such systems (cf.~III, (5.12)):
\begin{equation}
\sum_\mathbf{p}[\psi_{pi}(\mathbf{r})\psi^*_{pk}(\mathbf{r}') + \psi_{-pi}(\mathbf{r})\psi^*_{-pk}(\mathbf{r}')] = \delta_{ik}\delta(\mathbf{r} - \mathbf{r}').
\tag{75.5}
\end{equation}

\SourcePage{332}{337}
The sum on the right-hand side of~(75.4) differs from the sum written above by the replacement of $\psi^*_k$ by $(\psi^*\gamma^0)_k$, and equals $\gamma^0_{kk'}\delta(\mathbf{r} - \mathbf{r}')$. Thus,
\begin{equation}
\{\hat{\psi}_i(\mathbf{r},t),\hat{\bar{\psi}}_k(\mathbf{r}',t)\}_+ = \delta(\mathbf{r} - \mathbf{r}')\gamma^0_{ik}.
\tag{75.6}
\end{equation}

Let us note that from this formula it follows, in particular, the assertion mentioned in~\S\,74, about the anticommutativity of the operators $\hat{\psi}$ and $\hat{\bar{\psi}}$ outside the light cone. When $(x - x')^2 < 0$ there always exists such a system of reference in which $t = t'$; while at the same time $\mathbf{r} \neq \mathbf{r}'$, so that the anticommutator~(75.6) is effectively equal to zero.

Substituting~(75.6) into~(75.3) (and dropping the bispinor indices), we find finally\,\footnote{In the explicit notation with bispinor indices
\begin{equation}
(\gamma\hat{p} - m)_{ik}\,G_{kl}(x - x') = \delta^{(4)}(x - x')\delta_{il}.
\tag{75.7a}
\end{equation}}
\begin{equation}
(\gamma\hat{p} - m)\,G(x - x') = \delta^{(4)}(x - x').
\tag{75.7}
\end{equation}

Thus, the electron propagator satisfies the Dirac equation with a $\delta$-function on the right-hand side. In other words, it is the \textit{Green's function} of the Dirac equation.

In the following we shall have to deal not with the function $G(\xi)$ itself ($\xi = x - x'$), but with its Fourier components
\begin{equation}
G(p) = \int G(\xi)\,e^{ip\xi}\,d^4\xi
\tag{75.8}
\end{equation}
(the propagator in the momentum representation). Taking the Fourier component of both sides of~(75.7), we find that $G(p)$ satisfies the system of algebraic equations
\begin{equation}
(\gamma p - m)\,G(p) = 1.
\tag{75.9}
\end{equation}

The solution of this system:
\begin{equation}
G(p) = \frac{\gamma p + m}{p^2 - m^2}.
\tag{75.10}
\end{equation}

The four components of the 4-vector $p$ in $G(p)$ are independent variables (not connected with the relativistic equation $p^2 \equiv p_0^2 - \mathbf{p}^2 = m^2$). Writing the denominator in~(75.10) in the form $p_0^2 - (\mathbf{p}^2 + m^2)$, we see that $G(p)$ as a function of $p_0$ at a given $\mathbf{p}$ has two poles: at $p_0 = \pm\varepsilon$, where $\varepsilon = \sqrt{\mathbf{p}^2 + m^2}$. In the integration over $dp_0$ in the integral

\SourcePage{333}{338}
\begin{equation}
G(\xi) = \frac{1}{(2\pi)^4}\int e^{-ip\xi}\,G(p)\,d^4p = \frac{1}{(2\pi)^4}\int d^3p\cdot e^{i\mathbf{p}\mathbf{r}}\int dp_0\cdot e^{-ip_0\tau}\,G(p)
\tag{75.11}
\end{equation}
($\tau = t - t'$) the question therefore arises about the manner of going around the poles; without indicating this manner expression~(75.10) is essentially still indeterminate.

To elucidate this question let us return to the original definition~(75.1). Substituting into it the $\psi$-operators in the form of the sums~(73.6), and noting that from the vacuum average of zero only the following products of creation and annihilation operators are different from zero:
\[
\langle 0|a_\mathbf{p} a_\mathbf{p}^+|0\rangle = 1,\qquad \langle 0|b_\mathbf{p} b_\mathbf{p}^+|0\rangle = 1.
\]
(Since in the vacuum state there are no particles, before ``annihilating'' a particle by the operator $\hat{a}_\mathbf{p}$ or $\hat{b}_\mathbf{p}$, one must first ``create'' it by its operator $\hat{a}_\mathbf{p}^+$ or $\hat{b}_\mathbf{p}^+$.) We obtain
\begin{equation}
\begin{split}
G_{ik}(x - x') &= -i\sum_\mathbf{p}\psi_{pi}(\mathbf{r},t)\bar{\psi}_{pk}(\mathbf{r}',t) =\\
&= -i\sum_\mathbf{p}e^{-i\varepsilon(t-t')}\psi_{pi}(\mathbf{r})\bar{\psi}_{pk}(\mathbf{r}'),\quad t - t' > 0;
\end{split}
\tag{75.12}
\end{equation}
\[
G_{ik}(x - x') = i\sum_\mathbf{p}\bar{\psi}_{-pk}(\mathbf{r}',t')\psi_{-pi}(\mathbf{r},t) =
\]
\[
= i\sum_\mathbf{p}e^{i\varepsilon(t-t')}\psi_{-pi}(\mathbf{r})\bar{\psi}_{-pk}(\mathbf{r}')
\]
(when $t - t' > 0$ only the electronic terms contribute to $G$, and when $t - t' < 0$ only the positronic ones).

Imagining the summation over $\mathbf{p}$ to be replaced by integration over $d^3p$, and comparing~(75.12) with~(75.11), we see that the integral
\begin{equation}
\int e^{-ip_0\tau}G(p)\,dp_0
\tag{75.13}
\end{equation}
should have the phase factor $e^{-i\varepsilon\tau}$ when $\tau > 0$ and $e^{i\varepsilon\tau}$ when $\tau < 0$. We satisfy this condition if we agree to go around the poles $p_0 = \varepsilon$ and $p_0 = -\varepsilon$ respectively from above and from below (in the plane of the complex variable $p_0$):
% [Diagram (75.14): Complex $p_0$-plane with poles at $+\varepsilon$ (above real axis) and $-\varepsilon$ (below real axis), showing integration contour]

Indeed, when $\tau > 0$ we close the integration contour by an infinitely distant semicircle in the lower half-plane, so that the value of the integral~(75.13) will be given by the residue in the pole $p_0 = +\varepsilon$; when $\tau < 0$ we close the contour in the upper half-plane, and the integral is determined by the residue in the pole $p_0 = -\varepsilon$. In both cases the required result is obtained.

\SourcePage{334}{339}
This rule of going around (the \textit{Feynman rule}) can be formulated otherwise: integration is performed everywhere along the real axis itself, but to the mass of the particle there is ascribed an infinitely small negative imaginary part:
\begin{equation}
m \to m - i0.
\tag{75.15}
\end{equation}
Indeed, then
\[
\varepsilon \to \sqrt{\mathbf{p}^2 + (m - i0)^2} = \sqrt{\mathbf{p}^2 + m^2 - i0} = \varepsilon - i0.
\]
In other words, the poles $p_0 = \pm\varepsilon$ shift downward and upward from the real axis:
% [Diagram (75.16): Complex $p_0$-plane with displaced poles at $+\varepsilon - i0$ (below real axis) and $-\varepsilon + i0$ (above real axis)]
so that integration along the real axis becomes equivalent to integration along the contour~(75.14)\,\footnote{It is useful to note that the rule of pole displacement corresponds to the fact that $G(x - x')$ acquires an infinitely small damping at $|\tau|$, where $\tau = t - t'$. Indeed, if one writes the value of $p_0$ in the displaced poles as $-(\varepsilon - i\delta)$ and $+(\varepsilon - i\delta)$ (where $\delta \to +0$), then the temporal factor in the integral~(75.13) will equal $\exp(-i\varepsilon|\tau| - \delta|\tau|)$.}. Taking into account the rule~(75.15), the propagator~(75.10) can be written in the form
\begin{equation}
G(p) = \frac{\gamma p + m}{p^2 - m^2 + i0}.
\tag{75.17}
\end{equation}

The rule of integration under displacement of the poles is demonstrated by the following relation:
\begin{equation}
\frac{1}{x + i0} = \mathrm{P}\frac{1}{x} - i\pi\delta(x).
\tag{75.18}
\end{equation}
It must be understood in the sense that upon multiplication by any function $f(x)$ and integration we have
\begin{equation}
\int_{-\infty}^{\infty}\frac{f(x)\,dx}{x + i0} = \int_{-\infty}^{\infty}\frac{f(x)}{x}\,dx - i\pi f(0),
\tag{75.19}
\end{equation}
where the crossed-out integral sign, or the symbol P, denotes the principal value.

\SourcePage{335}{340}
The Green's function~(75.10) represents the product of the bispinor factor $\gamma p + m$ and a scalar:
\begin{equation}
G^{(0)}(p) = \frac{1}{p^2 - m^2}.
\tag{75.20}
\end{equation}
The corresponding coordinate function $G^{(0)}(\xi)$ is, evidently, a solution of the equation
\begin{equation}
(\hat{p}^2 - m^2)\,G^{(0)}(x - x') = \delta^{(4)}(x - x'),
\tag{75.21}
\end{equation}
i.e.\ the Green's function of the equation $(\hat{p}^2 - m^2)\psi = 0$. In this sense one can say that $G^{(0)}(x - x')$ is the propagator of scalar particles. It is easy to verify by computation (analogous to the one performed above) that the propagator of a scalar field is expressed through the $\psi$-operators~(11.2) by the formula
\begin{equation}
G^{(0)}(x - x') = -i\langle 0|\mathrm{T}\psi(x)\psi^+(x')|0\rangle.
\tag{75.22}
\end{equation}
analogous to definition~(75.1). In this the chronological product is defined (as for all bosonic operators) in the following manner:
\begin{equation}
\mathrm{T}\hat{\psi}(x)\hat{\psi}^+(x') = \begin{cases}\hat{\psi}(x)\hat{\psi}^+(x'),& t > t',\\[3pt] \hat{\psi}^+(x')\hat{\psi}(x),& t < t'\end{cases}
\tag{75.23}
\end{equation}
(with the same signs when $t > t'$ and $t < t'$).
